% 3-6-baselines.tex
%  Budget: 4pp.  Converted from thesis/drafts/3-6-baseline-models.md.
%  Source map: the "similar day" naive rule, the LEAR-LASSO reference model,
%   the 1092-day calibration window and the published forecast files used for
%   the end-to-end verification -> lago_forecasting_2021 ;
%   earlier LEAR/DNN treatment in the same benchmark line -> lago_forecasting_2018 ;
%   per-hour decomposition as standard practice in EPF -> jiang_review_2018 .
%  Numbers verified against reports/tables/results_canonical.csv and
%   lago_comparison.csv: naive 7.750/0.849 ; SARIMAX 4.351/0.477 ;
%   LEAR-LASSO 3.899/0.427 ; LEAR 1092 (shipped) 3.930 ; DNN Ensemble 3.413 ;
%   LEAR Ensemble (shipped) 3.609.
%  DROPPED from the draft: a year-by-year split (2016: 3.452 vs 3.474) that is
%   not present in any frozen table. Not quoted rather than quoted unverified.

\section{مدل‌های پایه}
\label{sec:3-6}

هر سه مدل این بخش بر یک واسط مشترک پیاده‌سازی شده‌اند و همگی از منبع یکتای ویژگیِ
\cref{sec:3-4} تغذیه می‌کنند. در نتیجه چارچوب اعتبارسنجیِ \cref{sec:3-5} هر سه را بدون هیچ حالت
خاصی می‌راند، و تفاوت عملکرد تنها به خودِ مدل بازمی‌گردد و نه به تفاوت در خوراک داده.

یکدستیِ یادشده از راحتیِ مهندسی نمی‌آید؛ تصمیمِ طراحیِ آگاهانه‌ای است. هر مدلی که مسیر
ویژگی‌سازیِ مخصوص خود را داشته باشد، سطح تازه‌ای برای نشت اطلاعات می‌گشاید که آزمون
\cref{sec:3-4} آن را پوشش نمی‌دهد؛ و هر تفاوتی در پنجره‌ی آموزش یا ترتیب ستون‌ها، بخشی از
اختلاف دقت را به عاملی جز خودِ مدل نسبت می‌دهد. سه مدل این بخش، مرجع‌هایی هستند که
مدل‌های \cref{sec:3-7} در برابر آن‌ها سنجیده می‌شوند؛ اعتبار آن سنجش تماماً به یکسان‌بودن
شرایط بستگی دارد.

\subsection*{مدل ساده‌ی مرجع}

مدل ساده‌ی «روز مشابه»، مطابق تعریف مرجع \cite{lago_forecasting_2021}، چنین است: برای
سه‌شنبه تا جمعه، قیمت‌های ۲۴گانه‌ی روز پیش (\lr{D-1})؛ برای دوشنبه، قیمت‌های جمعه‌ی پیش
(\lr{D-3})؛ و برای شنبه و یکشنبه، قیمت‌های همان روز در هفته‌ی پیش (\lr{D-7}).

منطق این قاعده رفتاری است و شاهد تجربی‌اش در \cref{sec:3-3-3} آمد: دوشنبه با یکشنبه‌ی پیش از
خود شباهت ساختاری ندارد، پس به آخرین روز کاری ارجاع داده می‌شود، و روزهای تعطیل به
هم‌تای خود در هفته‌ی پیشین.

این مدل هیچ پارامتری برازش نمی‌کند و صرفاً انتخاب ستون از ماتریس ویژگیِ مشترک است؛
متد برازش آن عامدانه بی‌اثر است و تنها برای پایبندی به واسط مشترک وجود دارد. لازمه‌ی
آن این است که فهرست وقفه‌های قیمتی در پیکربندی، دربردارنده‌ی وقفه‌های ۱، ۳ و ۷ باشد،
که پیکربندی پیش‌فرض آن را برآورده می‌کند.

نقش این مدل دوگانه است: کفِ عملکرد برای سنجش سایر مدل‌ها، و نقطه‌ی مرجعِ تفسیر معیار
\lr{rMAE}. بر دوره‌ی آزمون، \lr{MAE} ساعتی آن \lr{7.750} و \lr{rMAE} آن \lr{0.849}
است. درباره‌ی اینکه چرا این عدد برابر یک نیست، \cref{sec:3-5} توضیح داده است.

\subsection*{مدل \lr{SARIMAX}}

به‌عنوان نماینده‌ی رویکرد کلاسیک آماری، ۲۴ مدل \lr{SARIMAX} مستقل \lr{—} یکی به‌ازای
هر ساعت شبانه‌روز \lr{—} برازش می‌شود. این تجزیه‌ی ساعت‌به‌ساعت با ساختار مدل
\lr{LEAR-LASSO} هم‌قرینه است و با ادبیات بازار برق سازگار، زیرا هر ساعت از شبانه‌روز
رفتار آماری متمایزی دارد \cite{jiang_review_2018}؛ \cref{sec:3-3-3} دامنه‌ی این تفاوت را
نشان داد.

مرتبه‌ی مدل $(1,1,1)$ و مرتبه‌ی فصلی $(1,0,1)_{7}$ با دوره‌ی تناوب هفت‌روزه انتخاب شده
است \lr{—} یعنی فصلی‌بودن در مقیاس هفته و نه شبانه‌روز، چرا که بُعد ساعتی پیشاپیش با
تفکیک ۲۴گانه‌ی مدل‌ها جذب شده است. متغیرهای برون‌زا، پیش‌بینی‌های روزپیشِ بار و تولید
تجدیدپذیر برای خودِ روز هدف‌اند که، بنا بر استدلال \cref{sec:3-4}، پیش از مبدأ منتشر
می‌شوند و استفاده از آن‌ها نشت اطلاعات نیست.

برخلاف دو مدل دیگر این بخش، \lr{SARIMAX} افزوده‌ی خودِ این پژوهش است و در فهرست
مدل‌های مقاله‌ی مرجع جایی ندارد؛ نقش آن، تأمین یک مرجع کلاسیک و تفسیرپذیر در کنار
مرجع خطیِ ادبیات است.

\textbf{انحراف ثبت‌شده از پروتکل واسنجی روزانه.} برازش کامل ۲۴ مدل فصلی در ۷۲۸ مبدأ،
یعنی نزدیک به هفده هزار و پانصد برازش، عملی نیست؛ و برخلاف \lr{LEAR}، سابقه‌ای در
ادبیات مرجع برای آهنگ واسنجیِ \lr{SARIMAX} وجود ندارد که رعایت آن الزامی باشد. لذا
برازش کامل هر هفت روز انجام می‌شود و در میانه‌ی آن از به‌روزرسانیِ ارزانِ فضای حالت
استفاده می‌گردد: مشاهدات تازه به فیلتر افزوده می‌شوند و حالت مدل پیش می‌رود،
بی‌آنکه پارامترها دوباره بهینه شوند.

نکته‌ی مهم آنکه پیش‌بینی همچنان برای هر مبدأ روزانه تولید می‌شود \lr{—} گام چارچوب
پیش‌رونده همچنان یک روز است \lr{—} پس شمار و افق پیش‌بینی‌ها میان مدل‌ها یکسان و
مقایسه‌ی روزانه منصفانه باقی می‌ماند؛ تنها پارامترهای برازشیِ خودِ \lr{SARIMAX} است که
میان دو بازبرازش ثابت نگه داشته می‌شود. این تصمیم با تاریخ در دفتر تصمیمات ثبت شده
است، به‌جای آنکه در سکوت در کد بماند.

بر دوره‌ی آزمون، \lr{MAE} ساعتی این مدل \lr{4.351} و \lr{rMAE} آن \lr{0.477} است:
به‌روشنی بهتر از مدل ساده، اما ضعیف‌ترین عضو در میان مدل‌های برازشی.

\subsection*{مدل \lr{LEAR-LASSO}}

مدل \lr{LEAR}، یعنی رگرسیون خودتوضیح با تنظیم لاسو، مدل خطیِ مرجعِ ادبیات این حوزه
است \cite{lago_forecasting_2021, lago_forecasting_2018}. به‌ازای هر ساعت، یک رگرسیون
لاسو بر ۲۴۷ ویژگیِ \cref{sec:3-4} برازش می‌شود. تابع هدف لاسو برای بردار ضرایب $\beta$
چنین است:

\begin{equation}
\hat{\beta} = \arg\min_{\beta}
\left\{ \sum_{t=1}^{T}\left(p_t - x_t^{\top}\beta\right)^2
+ \lambda \sum_{j=1}^{J} \left|\beta_j\right| \right\}
\end{equation}

که در آن $p_t$ قیمت محقق‌شده در روز $t$، $x_t$ بردار ویژگی همان روز، $T$ شمار روزهای
پنجره‌ی واسنجی، $J$ شمار ویژگی‌ها و $\lambda$ ضریب جریمه است. جمله‌ی دوم، مجموع قدرمطلق
ضرایب را جریمه می‌کند و در نتیجه بسیاری از آن‌ها را دقیقاً صفر می‌کند؛ همین ویژگی
است که لاسو را برای ماتریسی با ۲۴۷ ستون و پنجره‌ای با ۱۰۹۲ سطر مناسب می‌سازد.

ضریب $\lambda$ با معیار اطلاعات آکائیکه انتخاب می‌شود و واسنجی مدل روزانه است.
پنجره‌ی واسنجی ۱۰۹۲ روز است، همان مقداری که در فایل پیکربندیِ ارزیابی به‌عنوان منبع
یکتای این عدد برای کل خط لوله تعریف شده است. نکته‌ی مهم آنکه $\lambda$ در هر برازش از
نو انتخاب می‌شود؛ یعنی این مدل ابرپارامترِ دستیِ تنظیم‌شده‌ای ندارد و برخلاف مدل‌های
\cref{sec:3-7} نیازمند جست‌وجوی \lr{Optuna} نیست.

پیاده‌سازی از واسط سطح‌پایینِ خودِ کتابخانه‌ی مرجع استفاده می‌کند، اما ماتریس‌های
ویژگی از خط لوله‌ی نشت‌آزموده‌ی \cref{sec:3-4} به آن داده می‌شود و نه از سازنده‌ی داخلی خودِ
کتابخانه. علت این انتخاب همان اصل ابتدای بخش است: اعتماد به یک سازنده‌ی ویژگیِ دوم،
مسیر دومی برای نشت می‌ساخت که آزمون‌های ما آن را نمی‌پوشاند. ترتیب و شمار ستون‌ها
دقیقاً با فرمول ویژگیِ مرجع منطبق است.

\textbf{راستی‌آزمایی زنجیره‌ی کامل.} صحت کل زنجیره \lr{—} بارگذار داده، مهندسی ویژگی،
چارچوب پیش‌رونده و پیاده‌سازی مدل \lr{—} با مقایسه‌ی مستقیم خروجی این مدل با
پیش‌بینی‌های خامِ منتشرشده‌ی خودِ مقاله‌ی مرجع بر همان دوره‌ی آزمون راستی‌آزمایی شد.

نخست هم‌سانیِ داده احراز شد: بیشینه‌ی اختلاف مطلق میان سری قیمت واقعیِ دو خط لوله، بر
تمام \lr{17{,}472} ساعت، صفر بود. سپس \lr{MAE} ساعتیِ مدل ما \lr{3.899} در برابر
\lr{3.930} برای گونه‌ی منتشرشده‌ی \lr{LEAR 1092} به دست آمد \lr{—} یعنی بازتولیدی وفادار
با همان پروتکل و همان داده.

\textbf{قیدی که باید همراه این عدد بیاید.} این هم‌ترازی را نباید به ادعای «شکست
بنچمارک» تعمیم داد. در همان مقایسه، گونه‌های ترکیبیِ منتشرشده همچنان پیش‌اند:
\lr{DNN Ensemble}\footnote{\lr{Deep Neural Network}} با \lr{3.413} و
\lr{LEAR Ensemble} با \lr{3.609}. آنچه این بخش
مجاز به ادعای آن است «کیفیت بازتولید» است، نه برتری. همین قید، مبنای تصمیم
ایست‌بازرسیِ میانیِ پروژه بود که خط دفاعی پژوهش را از شکستِ بنچمارک به مسیر
نوآوری‌محور تغییر داد. جدول کامل و آزمون‌های معناداریِ آن در \cref{sec:4-5-2} می‌آید.
